\chapter{Proposed Solution}\label{ch:proposed}

This chapter presents three mechanisms that target the structural
deficits of the three causal linear-time backbones formalised in
Section~\ref{sec:backbones}, together with the bidirectional
adaptation that extends each mechanism to the LION parallel form of
Section~\ref{sec:lion}. Section~\ref{sec:proposed:matrix} states the
problem formulation: the experimental matrix that the chapter is
organised around, the mechanism-prior alignment law that the matrix
is designed to test, and the probes used to test it.
Section~\ref{sec:proposed:msdc} introduces the Multi-Scale Depthwise
Convolution mechanism, which addresses the short-range temporal
hierarchy deficit identified in
Section~\ref{sec:backbones:deficit}. Section~\ref{sec:proposed:dho}
introduces the Damped Harmonic Oscillator mechanism, which extends
the real-diagonal transition operator of the diagonal-class backbones
with the block-complex transition primitive of
Section~\ref{sec:primitives:dho}. Section~\ref{sec:proposed:cvd}
introduces the Chunked Value Decorrelation mechanism, which targets
the absent interference-cancellation deficit on the linear-time
recurrent substrate. Section~\ref{sec:proposed:lion} describes the
unified LION wrapper that maps each causal backbone to its
bidirectional parallel form, and the decay-class correspondence that
follows. Section~\ref{sec:proposed:probes} closes the chapter with a
brief description of the diagnostic instrumentation used in
Chapter~\ref{ch:experiments} to attribute mechanism behaviour at the
parameter level.

\section{Problem formulation and the experimental matrix}\label{sec:proposed:matrix}

The three mechanisms introduced in the remainder of the chapter are
not the chapter's principal artefact. The principal artefact is the
\textit{transfer-pattern matrix} they populate. The mechanisms are
demonstrations of a predictive law about how a mechanism's gain on a
task relates to the structural deficit the architecture leaves
uncovered along the axis the mechanism targets; the matrix is the
structure on which that law becomes a falsifiable claim.

The experimental matrix is the three-by-two-by-two grid defined by
three causal backbones, two access modes, and two parameter scales.
The backbones are RWKV-6, Mamba-2, and Linear Attention, formalised
within the unified recurrence template of equation~(\ref{eq:unified})
in Sections~\ref{sec:backbones:rwkv6}, \ref{sec:backbones:mamba2},
and~\ref{sec:backbones:la} respectively. The two access modes are the
causal mode in which each backbone propagates state in one direction,
and the LION bidirectional mode of Section~\ref{sec:lion} in which
the parallel form gives every position access to the entire sequence.
The two parameter scales are seven million and fourteen million
parameters, corresponding to a six-layer and a twelve-layer encoder
respectively, with feature dimensions chosen so that the parameter
count is matched to within a small margin across backbones at each
scale. The three primary probes on which the matrix is evaluated are
LibriSpeech~\cite{panayotov2015librispeech} for the principal axis-one
comparison, Common Voice~\cite{ardila2019commonvoice} for
cross-distribution validation within axis one, and the Multi-Query
Associative Recall benchmark of~\citet{arora2023zoology} at sequence
lengths $T \in \{64, 256, 1024\}$ for axis two. The empirical content
of the matrix is reported in Chapter~\ref{ch:experiments}; this
chapter states the law the matrix is built to test and defines the
three mechanisms whose transfer it measures.

The predictive law that the matrix tests is the
\textit{mechanism-prior alignment} law: a mechanism's gain on a given
task is proportional to the residual structural deficit the
architecture leaves uncovered along the axis the mechanism targets.
The law is two-sided. On the positive side, a mechanism that extends
the function class along an axis on which the architecture leaves a
deficit, applied to a task that exercises that axis, should yield a
gain whose magnitude orders with the deficit size across
architectures. On the null side, a mechanism whose target axis is
not exercised by the task, or whose target axis the architecture's
native primitives already cover, should yield no measured gain even
if its parameters move under stochastic gradient descent. The matrix
is sized to support both kinds of prediction: each row exercises a
different combination of axis coverage and architectural primitive,
and the column structure permits the mode-conditional decomposition
of any observed transfer pattern.

The three mechanisms target distinct axes on different sides of
each backbone (input side, transition side, value side) and are
algebraically independent: cells in the experimental matrix
exercise individual mechanisms, pairwise compositions, and the
full triple composition.

\section{Multi-Scale Depthwise Convolution}\label{sec:proposed:msdc}

The first proposed mechanism is the Multi-Scale Depthwise Convolution
(MSDC). MSDC applies the input-side multi-scale local-mixing
pattern of Section~\ref{sec:primitives:msdc} to the linear-time
recurrent backbones of this thesis. The mechanism addresses the
short-range temporal hierarchy deficit identified in
Section~\ref{sec:backbones:deficit} as common to all three backbones:
none of the three has a native primitive that gives every token
explicit access to a short-range neighbourhood at the
phoneme-to-syllable scale that the ASR probe exercises.

\subsection{Generic form}

For an input sequence $x \in \mathbb{R}^{T \times D}$, MSDC applies a
parallel bank of depthwise dilated causal convolutions and sums the
results under per-channel learnable weights:
\begin{equation}\label{eq:msdc-generic}
\mathrm{MSDC}(x)[t] = \sum_{d \in \mathcal{D}} \alpha_{d} \odot \mathrm{DWConv}_{k, d}(x)[t],
\end{equation}
where $\mathcal{D} = \{1, 2, 4, 8\}$ is the dilation set, $k = 3$ is
the kernel size, $\mathrm{DWConv}_{k, d}$ denotes a depthwise
(channel-separable) one-dimensional causal convolution of kernel size
$k$ and dilation $d$, $\alpha_{d} \in \mathbb{R}^{D}$ are per-layer
learnable per-channel mixing weights, and $\odot$ denotes the
elementwise product over the channel dimension. The operator is
inserted on the input side of every backbone block: each block $f$
becomes $f(\mathrm{MSDC}(x))$, with the convolution kernels and the
mixing weights replicated independently per layer.

\subsection{Design notes}

The dilation set $\{1, 2, 4, 8\}$ at kernel size $k = 3$ produces a
maximum past-only receptive field of $1 + (k - 1) \cdot \max \mathcal{D} = 17$
frames at the deepest dilation. After the two-fold acoustic
subsampling that precedes the encoder in the experimental setup, the
range corresponds to roughly forty to one-hundred-sixty milliseconds
in the input audio, which matches the phoneme-to-syllable scale
relevant to character-level CTC decoding. The same dilation set was
selected on linguistic grounds in the autoregressive language model
of~\citet{kiruluta2025breaking} and is the most direct precedent for
the input-side multi-dilation form used here.

The depthwise structure of $\mathrm{DWConv}_{k, d}$ keeps the
parameter footprint of the operator at one kernel slot per channel
per dilation per layer, which is small relative to the time-mix
block on each backbone. The mixing weights $\alpha_{d}$ for $d > 1$
and the convolution kernels for $d > 1$ are initialised to a small
non-zero value rather than to zero, so that the gradient with respect
to $\alpha_{d}$ at step zero is non-degenerate. The choice follows
from the multiplicative form of equation~(\ref{eq:msdc-generic}): a
strict zero on either of the two factors at any single dilation
collapses the gradient through that dilation to zero, and the
operator silently reduces to its $d = 1$ branch.

\subsection{Per-backbone form}

The MSDC operator is architecture-agnostic on the input side, and
the three causal backbones receive identical insertion points modulo
the backbone's own time-mix entry interface.
Table~\ref{tab:msdc-deployment} contrasts the vanilla causal
recurrence of each backbone with its MSDC-augmented causal form and
with its MSDC-augmented LION form.

\begin{table}[t]
\centering
\small
\caption[Per-backbone form of Multi-Scale Depthwise
Convolution]{Per-backbone form of Multi-Scale Depthwise
Convolution. The vanilla causal column reproduces the backbone's
state and readout from Sections~\ref{sec:backbones:la},
\ref{sec:backbones:rwkv6}, and~\ref{sec:backbones:mamba2}. The two
augmented columns show the algebraic insertion point of MSDC on each
backbone in the causal mode and in the LION bidirectional mode. The
mechanism operates on the input side in every cell.}
\label{tab:msdc-deployment}
\setlength{\tabcolsep}{4pt}
\renewcommand{\arraystretch}{1.25}
\begin{tabularx}{\linewidth}{@{}p{2.1cm} >{\hsize=1.15\hsize\linewidth=\hsize}X >{\hsize=0.95\hsize\linewidth=\hsize}X >{\hsize=0.90\hsize\linewidth=\hsize}X@{}}
\toprule
\textbf{Architecture} & \textbf{Vanilla causal form} & \textbf{+ MSDC (causal)} & \textbf{+ MSDC (LION)} \\
\midrule
RWKV-6 &
$s_{t} = \mathrm{diag}(w_{t}) s_{t-1} + k_{t}^{\top} v_{t}$, see~(\ref{eq:rwkv6-state}), (\ref{eq:rwkv6-readout}) &
$\tilde{x} = \mathrm{MSDC}(x)$; the Time-Mix block consumes $\tilde{x}$ in place of $x$ and recomputes $w_{t}, k_{t}, v_{t}$ from $\tilde{x}$ &
LION-S parallel form on the MSDC-augmented Time-Mix; the WKV decay $w_{t}$ enters the LION mask of~(\ref{eq:lion-mask}) as the per-step factor $\lambda_{k}$ \\
\addlinespace
Mamba-2 &
\shortstack[l]{$h_{t} = a_{t} h_{t-1} + B_{t} x_{t}$\\see~(\ref{eq:mamba2-ssd})} &
$\tilde{x} = \mathrm{MSDC}(x)$; the SSD scan consumes $\tilde{x}$ in place of $x$ and recomputes $a_{t}, B_{t}, C_{t}$ from $\tilde{x}$ &
LION-S parallel form on the MSDC-augmented SSD path; the per-token scalar $a_{t}$ enters the LION mask as the per-step factor $\lambda_{k}$ \\
\addlinespace
\shortstack[l]{Linear\\Attention} &
$S_{t} = S_{t-1} + \phi(K_{t}) V_{t}^{\top}$, see~(\ref{eq:la-state}), (\ref{eq:la-output}) &
$\tilde{x} = \mathrm{MSDC}(x)$; the kernel state-update consumes $\tilde{x}$ in place of $x$ and recomputes $\phi(K_{t}), V_{t}, \phi(Q_{t})$ from $\tilde{x}$ &
LION-LIT parallel form on the MSDC-augmented kernel path; the controlled LION-S form of Section~\ref{sec:proposed:lion} additionally inherits a per-token sigmoid decay \\
\bottomrule
\end{tabularx}
\end{table}

The mechanism's algebraic point of action is identical across the
three rows: a pre-block input replacement
$x \leftarrow \mathrm{MSDC}(x)$ that the time-mix block of RWKV-6,
the SSD scan of Mamba-2, and the kernel state-update of Linear
Attention all read from. The table emphasises this consistency: the
augmented causal column differs across rows only in which
backbone-specific projections recompute their inputs from
$\tilde{x}$, and the augmented LION column differs across rows only
in the decay class the parallel form inherits from the underlying
causal backbone.

MSDC is logically of the same family as the input-side multi-scale
local-mixing pattern documented across vision, audio, and language
settings. Multi-dilation depthwise convolutions on the language side
appear in~\citet{kiruluta2025breaking}; depthwise two-dimensional
separable variants appear in
AudioRWKV~\cite{xiong2025audiorwkv}; learnable single-dilation
depthwise convolutions on the value path of a linear attention layer
appear in LiT~\cite{wang2025lit}; and the fixed quad-directional
shift of Vision-RWKV~\cite{duan2024visionrwkv} is a parameter-free
predecessor of the same family on two-dimensional vision input. The
MSDC variant in this thesis combines the multi-dilation form
of~\citet{kiruluta2025breaking} with the depthwise convolution
structure shared across the family, and applies it to three causal
linear-time backbones simultaneously.

\section{Damped Harmonic Oscillator}\label{sec:proposed:dho}

The second proposed mechanism is the Damped Harmonic Oscillator
(DHO). DHO applies the block-complex transition with Rayleigh
dissipation, introduced as a mathematical primitive in
Section~\ref{sec:primitives:dho}, to the diagonal-class backbones
whose vanilla transition is real-diagonal or trivial. The mechanism extends the function class of the
recurrence's transition operator from the real-diagonal class
$(\mathbb{R}_{+})^{K}$ to the smallest natural extension that admits
damped oscillation, the block-diagonal class
$SO(2)^{K/2} \times \mathbb{R}_{+}^{K/2}$. The standard reference for
the rotation group $SO(2)$ as a matrix Lie group is
~\citet[Chapter~1]{hall2015lie}; the standard reference for the
Rayleigh dissipation function from analytical mechanics
is~\citet[Sections~1.5 and~2.5]{goldstein2002classical}. The extended
transition remains inside the diagonal SSM class characterised
by~\citet{sarrof2024expressive}, since each two-by-two
rotation-and-damping block diagonalises over $\mathbb{C}$ to a pair
of conjugate eigenvalues $r e^{\pm i \theta}$ with $r = e^{-\lambda}$.

\subsection{Generic form}

DHO replaces the per-channel real-diagonal transition with a
block-diagonal product of $K/2$ rotation-and-damping blocks. Within
a single head $h$ and at step $t$, the block at index $b$ is
\begin{equation}\label{eq:dho-generic}
G_{t, h, b} = e^{-\lambda_{\mathrm{eff}, t, h, b}} \cdot R(\theta_{t, h, b}),
\qquad
\lambda_{\mathrm{eff}, t, h, b} = \lambda_{\mathrm{raw}, h, b} + \eta_{h, b} \cdot \theta_{t, h, b}^{2},
\end{equation}
where $R(\theta)$ is the SO(2) rotation matrix of
equation~(\ref{eq:block-complex}), $\theta_{t, h, b} \in \mathbb{R}$
is a data-dependent rotation angle produced from the input by a
low-rank LoRA-style projection, $\lambda_{\mathrm{raw}, h, b} \in \mathbb{R}$
is the per-head per-block raw damping log-rate, and
$\eta_{h, b} \ge 0$ is a per-head per-block viscosity coefficient
that couples the effective damping to the squared rotation angle
through the discrete Rayleigh-dissipation form of
equation~(\ref{eq:rayleigh}). The viscosity-coefficient parameter
shape is identical across all three backbones: a tensor of shape
$(n_{\mathrm{heads}}, K/2)$ per layer.

The rotation angle is bounded by a depth-graded clip schedule:
\begin{equation}\label{eq:dho-clip}
|\theta_{t, h, b, \ell}| \le \theta_{\max}(\ell),
\end{equation}
where $\ell$ is the layer index of the encoder and
$\theta_{\max}(\cdot)$ is a fixed depth-graded schedule that grows
from a tight clip in the early layers to a wider clip in the deeper
layers. For the six-layer encoder corresponding to the seven-million
parameter scale, the schedule takes the values
$(\pi/8, \pi/8, \pi/4, \pi/4, \pi/2, \pi/2)$ across layers one
through six. For the twelve-layer encoder corresponding to the
fourteen-million parameter scale, the schedule extends analogously,
with the same maximum clip preserved at the deepest layers. The
schedule is the mechanism's only canonical clip parameterisation.

\subsection{Design notes}

Three design choices specify the mechanism completely beyond the
algebraic form of equation~(\ref{eq:dho-generic}). First, the
viscosity coefficient $\eta_{h, b}$ is initialised to zero. At step
zero of training, the Rayleigh-dissipation term contributes nothing
to the effective damping,
$\lambda_{\mathrm{eff}, t, h, b} = \lambda_{\mathrm{raw}, h, b}$,
and DHO reduces exactly to the undamped block-complex transition of
equation~(\ref{eq:block-complex}). The mechanism is therefore
zero-regression-by-construction at initialisation: the augmented
recurrence at step zero is structurally identical to the same
backbone with a block-complex transition and no Rayleigh coupling,
and the additional capacity engages only as $\eta$ moves under the
gradient.

Second, the depth-graded $\theta$-clip schedule biases the rotation
operator toward near-identity initialisation in the shallow layers
and admits a wider rotation budget in the deeper layers. The
motivation is structural: shallow layers operate on input-aligned
representations for which a tight rotation budget keeps the
transition close to the diagonal class, while deeper layers operate
on more abstract phase structure for which a wider rotation budget
permits the block-complex extension to express richer dynamics.

Third, the data-dependent rotation angle $\theta_{t, h, b}$ is
produced by a low-rank LoRA-style projection of the input with a
depth-graded rank schedule that mirrors the depth-graded
$\theta$-clip schedule. For the six-layer encoder, the rank schedule
is $(16, 16, 32, 32, 48, 48)$ across layers one through six. For the
twelve-layer encoder, the schedule extends analogously. The pairing
of the two schedules ensures that the deeper layers, which receive
the wider rotation budget, also receive the higher-capacity LoRA
projection that can populate the budget.

Fourth, the per-block damping factor
$e^{-\lambda_{\mathrm{eff}, t, h, b}}$ in
equation~(\ref{eq:dho-generic}) makes the block-complex transition
decay-bounded by construction, independent of any external mask.
The property is local to the transition operator and does not
depend on the LION mask of equation~(\ref{eq:lion-mask}) supplying
its own decay through $\lambda_{k}$. DHO is therefore
decay-self-sufficient on either LION variant (LION-LIT or LION-S),
unlike the value-decorrelation mechanism of
Section~\ref{sec:proposed:cvd} whose conditioning guarantee
requires a decay-bearing mask.

\subsection{Per-backbone form}

DHO is a transition-side mechanism: it modifies the recurrent
operator $A_{t}$ of equation~(\ref{eq:unified}) and leaves the
input-write and the readout untouched.
Table~\ref{tab:dho-deployment} contrasts the vanilla transition of
each backbone with the DHO-augmented transition in the causal mode
and the DHO-augmented transition in the LION bidirectional mode.

\begin{table}[t]
\centering
\small
\caption[Per-backbone form of the Damped Harmonic
Oscillator]{Per-backbone form of the Damped Harmonic
Oscillator. The vanilla causal column reproduces the per-step
transition operator of each backbone from
Section~\ref{sec:backbones}. The augmented columns show the
DHO-replaced transition in the causal mode and in the LION
bidirectional mode. The block parameters $\eta_{h, b}$,
$\lambda_{\mathrm{raw}, h, b}$, and the LoRA projection that
produces $\theta_{t, h, b}$ are shared between the two modes per
backbone.}
\label{tab:dho-deployment}
\setlength{\tabcolsep}{4pt}
\renewcommand{\arraystretch}{1.25}
\begin{tabularx}{\linewidth}{@{}p{2.4cm} X X X@{}}
\toprule
\textbf{Architecture} & \textbf{Vanilla transition $A_{t}$} & \textbf{+ DHO (causal)} & \textbf{+ DHO (LION)} \\
\midrule
RWKV-6 &
$\mathrm{diag}(w_{t}) \in (\mathbb{R}_{+})^{D/h}$, see~(\ref{eq:rwkv6-decay}) &
block-diagonal product of $K/2$ blocks $G_{t, h, b}$ replacing pairs of real channels in $\mathrm{diag}(w_{t})$ &
LION-S parallel form whose mask factor inherits the block-complex transition $G_{t, h, b}$ in place of $w_{t}$ \\
\addlinespace
Mamba-2 &
$a_{t} I$ scalar identity, see~(\ref{eq:mamba2-ssd}) &
block-diagonal closure of pairs of state channels around the scalar $a_{t}$, with $\theta_{t, h, b}$ providing the per-block phase that $a_{t} I$ lacks &
LION-S parallel form whose mask factor inherits the block-complex closure of $a_{t} I$ \\
\addlinespace
\shortstack[l]{Linear\\Attention} &
$I$ identity, trivial &
block-diagonal product of $K/2$ rotation-and-damping blocks $G_{t, h, b}$ replacing the trivial transition &
controlled LION-S parallel form whose mask factor is the block-complex transition; see Section~\ref{sec:proposed:lion} \\
\bottomrule
\end{tabularx}
\end{table}

DHO has its greatest structural reach on the architecture whose
vanilla transition is most degenerate: Linear Attention, where the
identity transition admits no native attenuation and no native
phase, so the block-complex form introduces a function-class
capability that did not previously exist. The reach is intermediate
on Mamba-2, where the scalar-identity transition admits attenuation
but no phase; DHO contributes the phase. The reach is smallest on
RWKV-6, where the per-channel data-dependent decay vector $w_{t}$
already provides per-channel decay diversity, and DHO's contribution
is restricted to closing pairs of channels under rotation. The
deficit-proportional ordering of the table follows from the
structural reading: the more degenerate the vanilla transition, the
more residual function class the block-complex extension recovers.

The block-complex transition is structurally compatible with the
symmetric decay mask of the LION parallel form. The LION mask of
equation~(\ref{eq:lion-mask}) factorises as a product of per-step
decay factors $\lambda_{k}$, and the per-step factor $\lambda_{k}$
is replaced by the block-diagonal product of $K/2$
rotation-and-damping blocks under DHO, with no other change to the
parallel form. DHO in LION mode on RWKV-6 and Mamba-2 inherits
the LION-S decay class, since both backbones have a
data-dependent decay parameter that the LION mapping preserves.
On Linear Attention, the LION mode uses the controlled LION-S
configuration of Section~\ref{sec:proposed:lion}, in which a
per-token sigmoid decay is introduced as a separate structural
choice.

\section{Chunked Value Decorrelation}\label{sec:proposed:cvd}

The third proposed mechanism is the Chunked Value Decorrelation
(CVD). CVD is a value-side preconditioning operator constructed for
linear-time recurrent substrates. It addresses the absent
interference-cancellation deficit identified in
Section~\ref{sec:backbones:deficit}: when several stored keys are
similar, their values overlap in the bounded recurrent state and
become difficult to disentangle at readout, and none of the three
backbones has a native primitive that performs an active cancellation
of the correlated component before the value enters the state. The
mechanism shares the principle of key-similarity decorrelation with
LUCID Attention~\cite{duvvuri2026lucid} but is constructed for a
different substrate, and the relationship between the two is one of
shared logical motivation rather than direct adaptation.

\subsection{Generic form}

For each chunk of size $T_{c}$, with chunk-local keys
$K_{c} \in \mathbb{R}^{T_{c} \times d}$ and values
$V_{c} \in \mathbb{R}^{T_{c} \times d_{v}}$, CVD computes a
chunk-local preconditioner from key-similarity statistics and
applies its inverse to the values:
\begin{equation}\label{eq:cvd-generic}
\tilde{V}_{c} = P_{c}^{-1} V_{c},
\qquad
P_{c} = \exp\!\left( \tau_{h} \cdot \left[ \frac{K_{c}^{\mathrm{RN}} (K_{c}^{\mathrm{RN}})^{\top}}{\sqrt{d}} - \sqrt{d} \right] \right) + \varepsilon I,
\end{equation}
where $K_{c}^{\mathrm{RN}}$ denotes the row-RMS-normalised keys
within the chunk, $\tau_{h}$ is a per-head learnable temperature,
and $\varepsilon I$ is a small diagonal regulariser. The
preconditioner is the chunk-local analogue of the LUCID-faithful
preconditioner of equation~(\ref{eq:lucid-rn}), restricted to the
$T_{c} \times T_{c}$ block in place of the full $T \times T$
matrix of equations~(\ref{eq:lucid-base})
and~(\ref{eq:lucid-output}). The preconditioned values
$\tilde{V}_{c}$ are passed to the unmodified scan of the
underlying backbone in place of the raw values, so the
recurrence's transition operator and readout direction are not
changed.

The full $T \times T$ key-similarity matrix is never materialised.
The preconditioner is constructed and inverted on chunk-local
$T_{c} \times T_{c}$ blocks, which preserves the linear-time
backbone's $O(T)$ envelope at the expense of a $T_{c}^{3}$
per-chunk per-head linear-system solve cost, amortised across the
$T / T_{c}$ chunks of the sequence.

\subsection{Design notes}

Three design choices specify the mechanism beyond the algebraic
form of equation~(\ref{eq:cvd-generic}).

First, the per-head temperature is parameterised through a
softplus link
$\tau_{h} = \mathrm{softplus}(\tau_{\mathrm{raw}, h})$ to ensure
positivity. The initialisation of $\tau_{\mathrm{raw}}$ is chosen
so that $\tau_{h} = 1$ at the start of training on the RWKV-6,
Mamba-2, and LION-mode Linear Attention forms; the causal Linear
Attention form uses $\tau_{\mathrm{raw}} = 0$, which yields
$\tau_{h} = \ln 2$.

Second, the diagonal regulariser $\varepsilon$ takes two values
across the matrix conditional on backbone: $\varepsilon = 10^{-4}$
on Mamba-2 and $\varepsilon = 10^{-6}$ on RWKV-6 and Linear
Attention. The Mamba-2 value is larger because the SSD chunk-local
Gram matrix becomes ill-conditioned at trained temperatures around
$\tau_{h} \approx 1.5$; the smaller value on the other two
backbones is paired with an element-wise singular fallback in the
linear solve.

Third, the chunk size $T_{c} = 64$ is uniform across backbones,
matched to the SSD chunk size of Mamba-2 and to the WKV scan
chunk of RWKV-6. The LION parallel form preserves the chunk-local
preconditioner construction at the same $T_{c}$.

\subsection{Provenance: shared logic, distinct substrate}

The relationship between CVD and the LUCID Attention mechanism
of~\citet{duvvuri2026lucid} is one of shared logical motivation
rather than direct adaptation. LUCID is a softmax-attention
mechanism that operates on the full $T \times T$ attention matrix
inside a quadratic-complexity Transformer, preserves the
$O(N^{2} d)$ complexity envelope, constructs the preconditioner of
equation~(\ref{eq:lucid-base}) as a fixed function of the keys
without a learnable temperature, and combines the preconditioner
with the standard softmax weights of
equation~(\ref{eq:lucid-output}) in front of the value path. CVD
operates inside the recurrence of a linear-time RNN, never
materialises the full $T \times T$ matrix, uses chunk-local
$T_{c} \times T_{c}$ systems exclusively, introduces a per-head
learnable $\tau_{h}$ in place of LUCID's fixed $1 / \sqrt{d}$
scaling, and feeds the preconditioned values into the backbone's
unchanged scan rather than into a softmax-weighted attention sum.
The shared logical move is the use of a key-similarity
preconditioner to decorrelate values before they enter the memory
operation, motivated by the same observation that correlated keys
cause readout interference in any associative-memory
implementation. The relation to the rank-one Householder erasure
of DeltaNet~\cite{schlag2021deltanet, yang2024deltanet} is
parallel: both DeltaNet and CVD remove the component of each value
aligned with earlier keys under the key-similarity geometry, but
DeltaNet does so online with a single rank-one update per token
while CVD does so in closed form across the chunk-local block.

\subsection{Per-backbone form}

CVD has three backbone-specific forms because the chunk structure
of the substrate differs across backbones.
Table~\ref{tab:cvd-deployment} contrasts the chunk-local block,
the key-similarity source, and the substitution site on each
backbone. The mechanism is value-side: in every cell, it modifies
the values that enter the recurrence and leaves the recurrence's
transition operator and readout direction unchanged.

\begin{table}[t]
\centering
\small
\caption[Per-backbone form of Chunked Value
Decorrelation]{Per-backbone form of Chunked Value
Decorrelation. The mechanism operates on chunk-local
$T_{c} \times T_{c}$ blocks on every backbone. The
key-similarity source identifies the key-side tensor that
populates the chunk-local Gram matrix, and the substitution site
identifies where in the per-block computation the preconditioned
values $\tilde{V}_{c}$ replace the raw values.}
\label{tab:cvd-deployment}
\setlength{\tabcolsep}{4pt}
\renewcommand{\arraystretch}{1.25}
\begin{tabularx}{\linewidth}{@{}p{2.4cm} X X X@{}}
\toprule
\textbf{Architecture} & \textbf{Chunk-local block} & \textbf{Key-similarity source} & \textbf{Substitution site} \\
\midrule
RWKV-6 &
within-chunk inside the chunked WKV scan; $T_{c} = 64$ &
keys $K_{c}$ produced by the Time-Mix projection &
once per chunk per head, before the WKV scan consumes the values \\
\addlinespace
Mamba-2 &
within each chunk of the SSD dual form; $T_{c} = 64$, inherited from the SSD chunk size of~\citet{dao2024transformers} &
the chunk-local $C_{c}$ tensor, the query analog in the SSD scalar-identity dual form &
once per chunk per head, before the SSD scan consumes the values \\
\addlinespace
\shortstack[l]{Linear\\Attention} &
within-chunk in both modes; $T_{c} = 64$ inside the LION parallel form, and $T_{c} = 64$ inside the chunked causal kernel scan &
$\phi(K_{c})$ in the feature-mapped key sequence &
inside the parallel attention path in the LION mode; once per chunk in the causal mode \\
\bottomrule
\end{tabularx}
\end{table}

The Mamba-2 form uses the $C$-correlation analog rather than the
more obvious $B$-correlation analog. The choice is structural: in
the SSD scalar-identity dual form of
equation~(\ref{eq:mamba2-ssd}), the $C$ tensor is the query analog
and the $B$ tensor is the input-write analog, so the chunk-local
Gram matrix $C_{c} (C_{c})^{\top}$ corresponds directly to the
$K_{c}^{\mathrm{RN}} (K_{c}^{\mathrm{RN}})^{\top}$ Gram matrix of
the generic form and produces a preconditioner whose
interpretation matches the RWKV-6 and Linear Attention forms.

\section{Unified LION wrapper}\label{sec:proposed:lion}

The LION framework of~\citet{afzal2025linear}, presented in
Section~\ref{sec:lion}, supplies the bidirectional access pattern
for the experimental matrix on each of the three causal
backbones. This section describes the LION computation on each
backbone and the decay-class correspondence that maps each
causal backbone to its LION counterpart.

The LION computation is structured as a shared algebraic
convention rather than a single shared wrapper. Each causal
backbone in the matrix carries its own per-backbone form, and
the three forms share the algebraic structure of the
bidirectional parallel attention of equations~(\ref{eq:lion-mask})
and~(\ref{eq:lion-fullatt}) while differing in which backbone
parameter supplies the mask factor $\lambda_{k}$. On RWKV-6, the
LION mode reuses the Time-Mix block via a mode parameter,
dispatching into a parallel attention kernel that consumes the
WKV decay $w_{t}$ as the LION mask factor. On Mamba-2, the LION
mode reuses the SSD block, dispatching into a parallel attention
kernel that consumes the per-token scalar $a_{t}$ as the mask
factor. On Linear
Attention~\cite{katharopoulos2020transformers}, the LION form
runs through a separate kernel with its own parallel-form
forward path, since the underlying causal recurrence has no
decay parameter and the LION mask must be supplied independently.
The three forms present a uniform algebraic interface to the
rest of the chapter and to the experimental matrix, and the
chapter refers to them collectively as the LION wrapper.

The decay-class correspondence between each causal backbone and its
LION counterpart follows from the structure of the per-step decay
factor $\lambda_{k}$ in the LION mask of
equation~(\ref{eq:lion-mask}). RWKV-6 inherits LION-S, since the
WKV decay $w_{t}$ is per-channel data-dependent through
equation~(\ref{eq:rwkv6-decay}) and maps directly onto the LION-S
parameterisation of $\lambda_{k}$ as a learnable per-token gate.
Mamba-2 inherits LION-S, since the scalar $a_{t}$ is per-token
data-dependent through the selective step size $\Delta_{t}$ of
equation~(\ref{eq:mamba2-delta}) and maps onto the LION-S
parameterisation as a per-token, single-channel decay. Linear
Attention inherits LION-LIT, since the causal accumulator has
$A_{t} = I$ and therefore no decay to inherit, which corresponds
to the $\lambda_{k} = 1$ form of the LION mask.

The matrix additionally reports a controlled LION-S form on the
Linear Attention backbone, in which a per-token sigmoid decay is
introduced that is not present in the underlying causal backbone.
The controlled form serves as the cleanest available test of the
structural role of decay in the bidirectional setting: it
isolates the contribution of the per-token decay parameter from
every other aspect of the LION parallel form, since the Linear
Attention backbone has no other decay-bearing primitive, and it
allows the experimental matrix to compare LION-LIT and LION-S on
the same kernel-attention substrate. The structural motivation
for the controlled form is recalled from Section~\ref{sec:lion}:
bounded value accumulation in the bidirectional setting requires
either an external normalisation step or a decay-bearing mask,
and a downstream value-decorrelation mechanism such as the
Chunked Value Decorrelation of Section~\ref{sec:proposed:cvd}
loses its conditioning guarantee when applied to a substrate
whose accumulation is unbounded.
Table~\ref{tab:lion-mapping} summarises the per-backbone mapping.

\begin{table}[t]
\centering
\small
\caption[LION decay-class mapping per backbone]{Decay-class
correspondence between each causal backbone and its LION
counterpart. The default LION class is the natural mapping per the
LION framework of~\citet{afzal2025linear}. The controlled LION-S
form on Linear Attention introduces a per-token sigmoid decay
that is not present in the underlying causal backbone, and is
the only entry in the controlled column.}
\label{tab:lion-mapping}
\setlength{\tabcolsep}{4pt}
\renewcommand{\arraystretch}{1.25}
\begin{tabularx}{\linewidth}{@{}p{2.6cm} X X X@{}}
\toprule
\textbf{Architecture} & \textbf{Causal decay class} & \textbf{LION decay class (default)} & \textbf{LION decay class (controlled)} \\
\midrule
RWKV-6 & per-channel data-dependent & LION-S & not applicable \\
\addlinespace
Mamba-2 & scalar per-token data-dependent & LION-S & not applicable \\
\addlinespace
Linear Attention & none ($\lambda_{k} = 1$) & LION-LIT & LION-S \\
\bottomrule
\end{tabularx}
\end{table}

\section{Diagnostic instrumentation}\label{sec:proposed:probes}

The empirical chapter attributes mechanism behaviour at the
parameter level using a set of per-mechanism diagnostic probes
whose construction is summarised here. Each probe reports a
per-step scalar that distinguishes whether a mechanism is engaged
under stochastic gradient descent and whether its engagement is
helpful for the primary metric, supporting the four-cell
engagement classifier that organises the closed-cell evidence in
Chapter~\ref{ch:experiments}.

For the Multi-Scale Depthwise Convolution mechanism, the probe
reports the per-layer norm of the mixing weights $\alpha_{d}$ for
$d > 1$ and the per-layer gradient signal on those weights,
distinguishing the case in which the multi-dilation branches remain
at their initialisation from the case in which they diverge under
the gradient. For the Damped Harmonic Oscillator mechanism, the
probe reports the per-head per-block trajectory of $\eta_{h, b}$
and the rotation-angle distribution $\theta_{t, h, b}$,
distinguishing the case in which the Rayleigh-dissipation coupling
remains near zero from the case in which it engages. For the
Chunked Value Decorrelation mechanism, the probe reports the
trained per-head temperature $\tau_{h}$ and the chunk-local
condition number of the Gram matrix that enters $P_{c}$,
distinguishing the case in which the preconditioner is near-identity
from the case in which it actively decorrelates. The four-cell
engagement classifier crosses the parameter-mobility signal with
the primary metric outcome, separating mechanisms that are engaged
and helpful from those that are engaged without converting into
measured gain, those that fail to engage, and the residual cell.
The full classifier definition and its application to the
closed-cell evidence are deferred to Chapter~\ref{ch:experiments}.
